\documentclass[11pt]{article}
%\documentclass[a4paper,11pt,twocolumn]{article}%twocolumn layout; not sure if this works correctly. Feel free to experiment.
\usepackage[pdftex]{color,graphicx}
\usepackage[T1]{fontenc}
\usepackage{pxfonts}
% \usepackage{subfigure}
\usepackage{xcolor}

\usepackage{svg}

%color links to figs, etc.
%\usepackage[pdftex,colorlinks,urlcolor=blue,breaklinks]{hyperref}
\usepackage[hidelinks]{hyperref}
		    
%easily manipulate margins
\usepackage{geometry}

%CC logos to add in the footer
% \usepackage{ccicons}

\usepackage[utf8]{inputenc} % allow utf-8 input
\usepackage{url}            % simple URL typesetting
\usepackage{booktabs}       % professional-quality tables
\usepackage{amsfonts}       % blackboard math symbols
\usepackage{amssymb}
\usepackage{amsmath}
\usepackage{nicefrac}       % compact symbols for 1/2, etc.
\usepackage{microtype}      % microtypography
\usepackage{xcolor}         % colors
\usepackage{algorithmicx}
\usepackage{algorithm, algpseudocode}
\usepackage{graphicx}
\usepackage{adjustbox}
\usepackage{mathrsfs}

\usepackage{subcaption}



%enhanced floats
\usepackage{float}

%rotate floats if necessary
\usepackage{rotating}

%nicer, clever tables, uncomment if necessary
% \usepackage{supertabular}

%tables with notes, etc., uncomment if necessary
%\usepackage{threeparttable}

%improved captions, uncomment if necessary
\usepackage{caption}


%add nice headers
% \usepackage{fancyhdr}
% \pagestyle{fancy}
% \renewcommand{\headrulewidth}{0pt}
% \lhead{{\footnotesize }}
% \chead{}
% \rhead{{\footnotesize }}
% \lfoot{\ccbynd}
% \cfoot{\thepage}


%flushleft the title, authorblock and Abstract
%modified from http://tex.stackexchange.com/questions/85343/left-align-abstract-title-and-authors


\hyphenation{}

\renewcommand*{\thefootnote}{\fnsymbol{footnote}}

%if to do notes need to be added: need to configure this!
\usepackage[colorinlistoftodos]{todonotes}

%feel free to change the reference style to suit your needs
% \usepackage[firstinits=true, backref=false, maxcitenames=1,  hyperref=auto, style=authoryear, defernumbers=true, backend=bibtex]{biblatex}[2010/11-19]
\usepackage[style=numeric,sorting=nyt]{biblatex}

%change the name of this file to point to your bib file.
% \bibliography{./Bibliography/refs}
\addbibresource{./Bibliography/refs.bib} %Import the bibliography file

\newcommand{\yz}[1]{{\color{blue}[Yizi: #1]}}

\definecolor{IBMblue}{HTML}{648FFF}
\definecolor{IBMpurple}{HTML}{785EF0}
\definecolor{IBMhotpink}{HTML}{DC267F}
\definecolor{IBMorange}{HTML}{FE6100}
\definecolor{IBMyellow}{HTML}{FFB000}

\newcommand{\mup}{$\mu$\textbf{P}}




\begin{document}

\title{On Language Model Scaling Laws with SVG Code}
\date{Spring 2026}
\author{Vincent Hepola}


\maketitle
\tableofcontents




\section{Introduction}
% (0.5–1 page)

\subsection{Motivation}
% Motivation for studying scaling laws on SVG data

Why SVG?
\begin{itemize}
	\item Strict syntax rules, structure,  results immediately verifiable with visual inspection.
\end{itemize}

\subsection{Approach}
% Overview of your approach

Four steps
\begin{enumerate}
	\item Collect and preprocess SVG data, train a BPE tokenizer, ensure desired Token counts.
	\item Build the Transformer model architecture from scratch, short train 5 different model sizes at fixed learning rate (LR).
	\item Build new Transformer model with \mup, repeat short train of same 5 model sizes based on fixed learning rate.
	\item Full train of best model, generate samples.
\end{enumerate}





\section{Data}
% (1–1.5 pages)

\subsection{Dataset and Preprocessing}
% Dataset description and preprocessing pipeline

\subsection{Tokenization}
% Tokenization scheme with examples

\subsection{Statistics and Examples}
% Statistics and rendered SVG examples




\section{Methods}
\section{Results}
\section{Discussion}
\section{Conclusion}
\section{References}
\section{Appendices}




% \vspace{1em}


% \begin{figure}[htbp]
% 	\centering
% 	\begin{subfigure}[b]{0.6\textwidth}
% 		\centering
% 		\includegraphics[width=\textwidth]{figures/original LF.png}
% 		\label{fig:og}
% 	\end{subfigure}
% 	\begin{subfigure}[b]{0.6\textwidth}
% 		\centering
% 		\includegraphics[width=\textwidth]{figures/trilinear.png}
% 		\label{fig:tri}
% 	\end{subfigure}
% 	\begin{subfigure}[b]{0.6\textwidth}
% 		\centering
% 		\includegraphics[width=\textwidth]{figures/target HF.png}
% 		\label{fig:target}
% 	\end{subfigure}
% 	\caption{Original low-field, trilinear upsampled, and target high-field}
% 	\label{fig:triple}
% \end{figure}


% \begin{figure}[htbp]
% 	\centering
% 	\includegraphics[width = \textwidth]{figures/UNet3D_DrawIO.png}
% 	\caption{\textbf{Our 3D U-Net Model} }
% 	\label{fig:Model}
% \end{figure}

% \textit{Show your results on the outlined training/validation approach on the public dataset (including key metrics such as standard deviation, etc. where relevant).}


% \begin{figure}
% 	\centering

% 	\begin{tabular}{|c|c|c|}
% 		\hline
% 		Kaggle Team         & Public Score       & Private Score      \\
% 		\hline
% 		JayZ (1st place)    & 0.6538             & 0.6544             \\
% 		sota\_diffusion     & 0.6399             & 0.6343             \\
% 		The Synapses (Ours) & \underline{0.5404} & \underline{0.5454} \\
% 		baseline            & 0.5129             & 0.5152             \\
% 		\hline
% 	\end{tabular}

% 	\caption{Kaggle Results, higher is better. Our scores are \underline{underlined}.}
% \end{figure}

% \begin{figure}[htbp]
% 	\centering
% 	% First subfigure: taking up 45% of the text width
% 	\begin{subfigure}[b]{0.45\textwidth}
% 		\centering
% 		\includegraphics[width=\textwidth]{figures/nifti LF test.png}
% 		\caption{Test sample\_019 low-field}
% 		\label{fig:sub1}
% 	\end{subfigure}
% 	\hfill % This pushes the subfigures apart and fixes underfull \hbox warnings
% 	% Second subfigure: taking up 45% of the text width
% 	\begin{subfigure}[b]{0.45\textwidth}
% 		\centering
% 		\includegraphics[width=\textwidth]{figures/nifti model pred.png}
% 		\caption{Test sample\_019 super-resolution result}
% 		\label{fig:sub2}
% 	\end{subfigure}

% 	\caption{Comparison between low-field sample\_019 and super-resolution.}
% 	\label{fig:LF_Pred_Comparison}
% \end{figure}


% Please make sure to include references for what you cite. \textbf{LLMs often get these wrong, so please double check if you use LLMs here!}


% \printbibliography[heading=none]
%%%%%%%%%%%%%%%%%%%%%%%%%%%%%%%%%%%%%%%%%%%%%%%%%%%%%%%%%%%%


\end{document}